\section{Methodology: PEAR}
\label{sec:method}

PEAR (Patch-Embedding Activation Routing) is a training-free, closed-form ensembling framework designed to dynamically blend specialized expert adapters on-the-fly. Its core philosophy is grounded in simplicity and system elegance: by shifting the routing operation to the base model's first structural projection layer---the frozen Patch Embedding layer (Layer 0)---PEAR completely resolves both the Routing Paradox and the Early-Feature Loss Trade-Off. 

The complete architecture of PEAR is illustrated in the sections below, detailing how an incoming query stream is routed and ensembled across 100\% of the transformer blocks in a single, parallel forward pass.

\subsection{Early Representation Extraction (Layer 0)}
Let $X = \{x_1, \dots, x_B\}$ represent an incoming batch of $B$ query images, where the stream can be highly heterogeneous (containing samples from multiple distinct task domains). For each individual image $x_b$ in the batch, we project the raw pixel inputs through the base model's pre-trained, frozen Patch Embedding layer (Layer 0):
\begin{equation}
    Z_b = \text{PatchEmbed}(x_b) \in \mathbb{R}^{N \times D}
\end{equation}
where $N$ is the number of patch tokens, and $D$ is the intermediate representation dimension of the Vision Transformer (e.g., $D = 192$ for ViT-Tiny). 

To extract a compact, translation-invariant global early representation vector $z_b$ that captures the semantic essence of the image, we take the spatial average over the patch token dimension:
\begin{equation}
    z_b = \frac{1}{N} \sum_{i=1}^N Z_{b, i, :} \in \mathbb{R}^D
\end{equation}
This average-pooled vector $z_b$ serves as a high-fidelity, zero-cost representation of the task-specific manifold, obtained immediately at the boundary of the model's forward execution.

\subsection{Addressing the Global-Average-Color Routing Paradox and System Scaling}
A key technical consideration of Layer 0 routing in Vision Transformers is that the Patch Embedding layer is a linear projection of raw input pixels. Consequently, spatially average-pooling its tokens $z_b = \frac{1}{N} \sum_{i=1}^N Z_{b, i, :}$ is mathematically equivalent to taking a linear projection of the global average color and brightness of the input image:
\begin{equation}
    z_b \approx \text{PatchEmbed}\left(\text{AveragePool}(X_b)\right)
\end{equation}
In highly semantic or fine-grained visual classification tasks within the same visual domain, global average color lacks sufficient discriminative signal. While our high-fidelity synthetic sandbox simulates distinct subspace task coordinates to validate PEAR's ensembling calibration, in real-world deployments on complex semantic datasets, this "Global-Average-Color Routing Paradox" can cause representation bleed.

To resolve this limitation without running the heavy backbone twice, we outline three minimalist, low-overhead solutions:
\begin{enumerate}
    \item \textbf{Lightweight Pre-backbone Classifier:} Instead of average-pooling Layer 0, the system can deploy a tiny, ultra-fast pre-trained parametric classifier (e.g., a 3-layer CNN or MobileNet-v3) trained on raw pixels. This adds negligible computational overhead ($<1\%$ of the ViT backbone) but guarantees highly discriminative semantic task boundaries.
    \item \textbf{Attention-based Spatial Pooling:} Rather than simple average-pooling (which discards all spatial layout), spatial or attention-based pooling mechanisms can be evaluated on Layer 0 patch tokens to capture structural task-specific visual signatures.
    \item \textbf{Early-Layer Routing Compromise:} If tasks share identical low-level statistics but diverge semantically, the routing boundary can be shifted slightly deeper (e.g., to Layer 1 or Layer 2). Computing similarities at Layer 2 still bypasses the Routing Paradox for the vast majority of the network (Blocks 3--12) and incurs only a negligible late-adaptation penalty, proving that PEAR can adapt to any semantic level.
\end{enumerate}
While the Early-Layer Routing Compromise shifts the routing boundary slightly deeper (to layer $l_{\text{route}} \geq 1$) to leverage semantically richer activations, executing early blocks $l < l_{\text{route}}$ using unadapted weights while fine-tuned experts were trained with active early adapters introduces a subtle training-serving representational discrepancy at block $l_{\text{route}}-1$. Writing $h_b$ and $W_{\text{base}}$ as shorthand for $h_b^{(l_{\text{route}}-1)}$ and $W_{\text{base}}^{(l_{\text{route}}-1)}$, this is formulated as:
\begin{align}
    h_{b,\text{serving}}^{(l_{\text{route}})} &= \text{Block}_{l_{\text{route}}-1}\left(h_b; W_{\text{base}}\right) \nonumber \\
    h_{b,\text{ideal}}^{(l_{\text{route}})} &= \text{Block}_{l_{\text{route}}-1}\left(h_b; W_{\text{base}} + \Delta W_k\right)
\end{align}
where $\Delta W_k$ represents the trained expert LoRA weight delta $\Delta W_k^{(l_{\text{route}}-1)}$. At inference time, bypassing early-layer adapters of the chosen expert $k$ (since the true expert identity is only known after computing similarity on $h_b^{(l_{\text{route}})}$) results in a representational mismatch at the boundary.

To completely eliminate this boundary discrepancy, we formulate \textbf{Early-Layer Freezing during Training} (ELFT). Under ELFT, if the specialized adapters are destined for a multi-task serving configuration with a deeper routing boundary $l_{\text{route}}$, the operator simply freezes the early blocks $l < l_{\text{route}}$ during fine-tuning. This mathematically ensures that the trained weight deltas in these early blocks are zero:
\begin{equation}
    \Delta W_k^{(l)} = 0 \quad \forall l < l_{\text{route}}, \quad \forall k \in \{1, \dots, K\}
\end{equation}
By aligning the training-time network architecture with the serving-time inference path, ELFT completely neutralizes the boundary representational mismatch while preserving PEAR's full-depth ensembling benefits for the remaining blocks $l \geq l_{\text{route}}$.

While these remedies successfully resolve domain-level representation bleed, we explicitly qualify the scope of early-layer routing under conditions of extreme intra-domain semantic overlap. For fine-grained, single-domain serving scenarios where tasks share identical background manifolds but diverge only on highly specific, fine-grained semantic features (e.g., distinguishing fine-grained bird species from fine-grained dog breeds), early-layer representations may exhibit complete representational overlap. Under such extreme regimes, simple zero-shot patch centroids computed at early layers may suffer from representation bleed, potentially requiring routing to be conducted at deeper blocks where semantic classes are fully separated (re-introducing sequential latency delay). To maintain optimal $O(1)$ serving latency under fine-grained overlap, operators can deploy more advanced, offline closed-form alignment strategies. Specifically, using Centered Kernel Alignment (CKA) or orthogonal Procrustes projection, one can align early-layer representations with final head spaces over the calibration split, projecting early features into a highly discriminative latent space prior to similarity evaluation.

\subsection{Zero-Shot Patch Centroids (ZPC)}
To establish reference coordinates for each specialized task domain in the Layer 0 space, PEAR directly leverages class-wise **Zero-Shot Patch Centroids** (ZPC) computed offline from a tiny, weakly-supervised calibration set. Let $\mathcal{C}_{k, c}$ represent the calibration samples belonging to class $c \in \{1, \dots, C\}$ of task $k$. The ZPC centroid $\mu_{k, c} \in \mathbb{R}^D$ is defined as the mean Layer 0pooled representation:
\begin{equation}
    \mu_{k, c} = \frac{1}{|\mathcal{C}_{k, c}|} \sum_{s \in \mathcal{C}_{k, c}} z_s^{(0)} \in \mathbb{R}^D
\end{equation}
where $z_s^{(0)}$ is the Layer 0 pooled representation of calibration sample $s$. 

Utilizing these early-layer centroids ensures that both the queries and the reference anchors live in the exact same Layer 0 representational manifold, completely resolving any cross-space representational misalignment. Because these centroids are computed offline as the simple mean of a tiny calibration set ($|\mathcal{C}_{k}| = 64$), this step introduces **zero new trainable parameters**, preserving our strictly parameter-free ensembling guarantees.

However, from a computational complexity perspective, evaluating similarity against class-wise centroids scales as $\mathcal{O}(K \times C)$, where $K$ is the number of expert adapters and $C$ is the number of classes per task. While this is mathematically instantaneous for small-to-medium class spaces (such as the 10 classes of CIFAR-10), deploying PEAR on large-vocabulary tasks (e.g., ImageNet with 1000 classes) could introduce a slight memory and distance computation bottleneck on highly resource-constrained edge NPUs. To mitigate this scaling cost, edge operators can adopt two systems-aware optimizations: (i) \textbf{Task-Level Centroid Routing}, where class-specific centroids are average-pooled into a single global task centroid $\mu_k \in \mathbb{R}^D$, reducing complexity to $\mathcal{O}(K)$ distance evaluations, or (ii) \textbf{Hierarchical Anchoring}, where class-specific centroids are offline-clustered (using K-Means) into a small, fixed set of $C^* \ll C$ representative cluster centroids to preserve intra-task representational shape while guaranteeing a strict upper bound on test-time similarity computations.

\subsection{Cosine Similarity on the Unit Hypersphere (Unit-Norm Projection)}
At inference time, the system must evaluate the similarity of the incoming global early representation $z_b$ against the ZPC centroids. Rather than using Euclidean distance (which is sensitive to representation magnitude shifts across layers and inputs), PEAR restricts similarities to a unit-norm sphere.

The raw similarity $s_{k, b}$ of query representation $z_b$ with respect to task expert $k$ is formulated by taking the maximum cosine similarity over the expert's class centroids:
\begin{equation}
    s_{k, b} = \max_{j} \text{cos\_sim}(z_b, \mu_{k, j}) = \max_{j} \frac{z_b \cdot \mu_{k, j}}{\|z_b\|_2 \|\mu_{k, j}\|_2}
\end{equation}
By projectively mapping representations to a unit-norm hypersphere and evaluating fine-grained similarity over class centroids, we achieve strict scale invariance, neutralizing representation drift and magnitude biases.

\subsection{Intra-Task Dispersion Calibration}
While Unit-Norm Calibration provides scale invariance, distinct task manifolds naturally exhibit different representational densities and intra-task variances. A tighter, more concentrated task manifold will naturally yield higher cosine similarities to its anchors compared to a highly dispersed, heterogeneous task manifold. 

To resolve this asymmetry, PEAR introduces Intra-Task Dispersion Calibration (IDC). During the offline calibration phase using the calibration set $\mathcal{C}_k$ of $|\mathcal{C}_k| = 64$ samples per task, we compute the expected in-distribution similarity scale $d_k$ of the calibration samples to their respective class centroids:
\begin{equation}
    \label{eq:idc}
    d_k = \frac{1}{|\mathcal{C}_k|} \sum_{s \in \mathcal{C}_k} \text{cos\_sim}(z_s, \mu_{k, c_s})
\end{equation}
where $c_s \in \{1, \dots, C\}$ is the class label of calibration sample $s$.

At test time, the raw cosine similarity $s_{k, b}$ is calibrated by dividing it by the expected dispersion factor $d_k$, yielding the dispersion-calibrated similarity $u_{k, b}$:
\begin{equation}
    u_{k, b} = \frac{s_{k, b}}{d_k}
\end{equation}
This calibration standardizes the similarity scales across diverse task manifolds, ensuring that more diverse tasks are not systematically outvoted by more concentrated ones.

\subsection{Temperature-Scaled Softmax \& OOD Rejection}
The calibrated similarities are normalized into sample-specific routing weights $\alpha_{k, b}$ using a temperature-scaled Softmax:
\begin{equation}
    \alpha_{k, b} = \frac{\exp(u_{k, b} / \tau)}{\sum_{j=1}^K \exp(u_{j, b} / \tau)}
\end{equation}
where $\tau > 0$ is a temperature hyperparameter that controls the sharpness of the routing decision. Smaller values of $\tau$ yield crisper, near-categorical routing, while larger values allow soft ensembling. In our empirical evaluations, we employ $\tau = 0.05$ as our default ensembling temperature under standard overlapping and non-linear sandbox layouts (where softer temperature acts as a strong representational regularizer and denoiser). In highly optimized expert regimes (low-noise settings), we utilize a sharp temperature of $\tau = 0.001$ to perform crisp, near-discrete routing.

To protect the system against Out-of-Distribution (OOD) inputs or queries that do not match any of the specialized task manifolds, PEAR incorporates a hard security threshold $\gamma_{\text{OOD}}$. If the maximum raw similarity is below this threshold, the routing weights default to a uniform distribution, executing a **Static Uniform Weight Merging Fallback** over all expert adapters:
\begin{equation}
    \alpha_{k, b} = \frac{1}{K} \quad \forall k \quad \text{if } \max_{j} s_{j, b} < \gamma_{\text{OOD}}
\end{equation}
This OOD fallback strategy is mathematically robust and avoids predictable prediction biases: rather than zeroing out activations (which would nullify logits and bias predictions toward Class 0 in a multi-head classifier setup), falling back to a uniform ensembling of experts ($\alpha_{k, b} = 1/K$) ensures that the model preserves representational diversity and performs robust general-purpose classification under OOD query streams.

While a static global threshold $\gamma_{\text{OOD}}$ is straightforward, different task manifolds naturally exhibit vastly different representational densities and intrinsic dispersions. For highly dispersed, high-noise task manifolds (such as SVHN), their raw cosine similarities are systematically lower than those of highly concentrated, low-noise manifolds (such as MNIST). Consequently, a static global threshold secure enough to protect MNIST will systematically over-reject valid queries on noisier tasks, collapsing their accuracies.

To resolve this density asymmetry without manual hand-tuning, we formulate \textbf{Adaptive Task-Specific Thresholding}:
\begin{equation}
    \gamma_{\text{OOD}, k} = \eta \cdot d_k
\end{equation}
where $\eta \in [0, 1]$ is a user-defined global security factor, and $d_k$ is the Intra-Task Dispersion Calibration factor computed in Equation~\ref{eq:idc} during the offline calibration phase.

At test time, the OOD rejection rule becomes task-adaptive: a query $x_b$ is rejected and defaults to the Static Uniform Weight Merging Fallback (or Hard Edge Rejection) if and only if:
\begin{equation}
    s_{k, b} < \gamma_{\text{OOD}, k} \quad \forall k \in \{1, \dots, K\}
\end{equation}
By scaling the security thresholds dynamically with each task's expected representational density $d_k$, this adaptive mechanism automatically lowers the rejection barrier for noisier, highly dispersed manifolds while maintaining strict, secure boundaries for concentrated ones. This resolves the security-selectivity trade-off elegantly and eliminates the need for manual hyperparameter hand-tuning under varying task distributions.

However, from a systems serving perspective, executing all $K$ specialized expert adapters simultaneously under this Static Uniform Weight Merging Fallback carries a distinct latency and memory overhead trade-off. It represents the maximum possible memory and computational bandwidth footprint at inference time, as all $K$ parallel adapter paths must be loaded and evaluated. On highly resource-constrained edge hardware where memory bandwidth is the primary bottleneck, this fallback would degrade serving latency to $O(K)$ if parallel hardware streams are exhausted.

To resolve this, edge operators can employ a \textbf{Hard Edge Rejection} fallback: setting $\alpha_{k, b} = 0 \quad \forall k$ and routing the query solely through the frozen base model, or routing through a single pre-selected generalist adapter. To avoid prediction logit nullification (which would trivially bias predictions toward Class 0 when setting adapter scales to zero under a multi-head expert classifier setup), the system employs a dedicated, task-agnostic, low-cost generalist classification head. This head consists of a lightweight, single-layer linear projection trained directly on the shared, unadapted base backbone representations. In real-world deployment scenarios, this generalist classification head is lightweightly optimized on the combined, weakly-supervised calibration splits ($K \times B_{\text{cal}} = K \times 64$ samples) for 5--10 epochs on CPU, requiring negligible computational overhead ($<1$ second of execution) and no external datasets. During a hard rejection event, the routing framework directly routes the backbone's final representation to this generalist head, completely bypassing task-specific expert heads. This guarantees that prediction logits are computed from a valid general-purpose distribution rather than being zeroed out, providing edge operators with a mathematically sound and extremely frugal fallback option.

\subsection{Dynamic Activation Blending (All Layers)}
Because the sample-specific ensembling weights $\alpha_{k, b}$ are computed entirely at Layer 0 (the Patch Embedding layer), they are immediately available before propagation through any subsequent transformer blocks. This enables PEAR to perform dynamic, sample-specific activation ensembling across \textbf{all subsequent layers} $l \in \{1, \dots, L\}$ of the network.

For each transformer block $l$, the output activation $h_b^{(l)}$ is computed on-the-fly as:
\begin{equation}
    h_b^{(l)} = h_b^{(l-1)} W_{\text{base}}^{(l)} + \sum_{k=1}^K \alpha_{k, b} \left( h_b^{(l-1)} A_k^{(l)} B_k^{(l)} \right)
\end{equation}
where $W_{\text{base}}^{(l)}$ represents the pre-trained, frozen base model weights, and $A_k^{(l)} \in \mathbb{R}^{D \times r}$ and $B_k^{(l)} \in \mathbb{R}^{r \times D}$ are the specialized, rank $r$ down- and up-projection matrices of expert LoRA adapter $k$. 

While Equation~13 formulates PEAR's ensembling over standard multi-head attention projection adapters, PEAR's activation-blending mechanics scale seamlessly to MLP (feed-forward) layers. For an MLP adapter consisting of down- and up-projections $A_{k,\text{mlp}}^{(l)}$ and $B_{k,\text{mlp}}^{(l)}$, the activation blending is computed identically as: $h_{b,\text{mlp}}^{(l)} = \text{MLP}_{\text{base}}(h) + \sum_{k=1}^K \alpha_{k,b} \left(\sigma(h A_{k,\text{mlp}}^{(l)}) B_{k,\text{mlp}}^{(l)}\right)$, where $\sigma$ is the non-linear activation. Crucially, because MLP layers typically contain over $60\%$ of the model's parameters and capture significant task-specific knowledge, adapting MLP adapters under PEAR is expected to yield even higher multi-task capacity. However, because MLP adapters are computationally wider than attention adapters, doing so increases the memory bandwidth footprint and thread-concurrency requirements, further magnifying the importance of our hardware scalability analyses and the Hard Edge Rejection fallback.

This formulation yields several significant system advantages:
\begin{itemize}
    \item \textbf{Single-Pass Execution:} The base model weights $W_{\text{base}}^{(l)}$ are executed exactly once per forward pass. The expert activations are scaled sample-wise and added dynamically, eliminating redundant backbone passes.
    \item \textbf{Flat $O(1)$ Latency:} The sequential depth complexity remains constant regardless of the number of experts $K$, as the parallel LoRA updates can be evaluated concurrently in parallel GPU streams.
    \item \textbf{No Dynamic Memory Buffers:} Since ensembling weights $\alpha_{k, b}$ are fixed at Layer 0, the system does not need to store intermediate activations to compute routing decisions at later layers, which dramatically reduces peak memory usage on edge devices.
    \item \textbf{Hardware and Memory Bandwidth Scalability Limits:} While sequential dependency depth latency is flat $O(1)$, loading and executing $K$ parallel expert pathways concurrently introduces an $O(K)$ computational (FLOPs) and memory bandwidth footprint. On highly resource-constrained edge NPUs or mobile devices with narrow memory bus widths, concurrent adapter executions are bounded by hardware memory transfer and thread concurrency ceilings. Specifically, if the memory bus cannot accommodate simultaneous loads of all $K$ expert adapters, physical memory transfer serialization occurs, degrading actual serving speeds and rendering the flat $O(1)$ latency assumption invalid under massive expert counts $K$. In such resource-constrained scaling regimes, the \textbf{Hard Edge Rejection} fallback (Section 3.6) serves as a vital systems-aware countermeasure, completely shutting off adapter computation and relying on the generalist head to bypass parallel adapter compute.
\end{itemize}
By combining these non-parametric mechanisms, PEAR represents the ultimate minimalist ensembling framework: high-performance dynamic ensembling achieved with zero trainable parameters and absolute system simplicity.
